\documentclass{article}
\usepackage[utf8]{inputenc}
\usepackage[T1]{fontenc}
\usepackage{geometry}
\geometry{
    a4paper,
    left=2.5cm,
    right=2.5cm,
    top=2.5cm,
    bottom=2.5cm
}
\usepackage{graphicx}
\usepackage{float}
\usepackage{hyperref}
\usepackage[english]{babel}
\usepackage{amsmath}
\usepackage{amssymb}
\usepackage{array}
\usepackage{booktabs}
\usepackage{caption}
\usepackage{subcaption}
\usepackage{enumitem}
\usepackage{fancyhdr}
\usepackage{titlesec}
\usepackage{url}
\usepackage{textcomp}
\usepackage{xcolor}
\usepackage{listings}
\usepackage{multirow}
\usepackage{longtable}
\usepackage{tabularx}
\usepackage{times}

% Set up page style
\pagestyle{fancy}
\fancyhf{}
\rfoot{\thepage}
\renewcommand{\headrulewidth}{0pt}

% Title formatting
\titleformat{\chapter}[display]{\normalfont\bfseries}{\chaptertitle\thechapter}{0pt}{\Huge}
\titlespacing*{\chapter}{0pt}{50pt}{40pt}

% Document info
\title{Template Document}
\author{Academic Support Agent}
\date{\today}

\begin{document}
\maketitle

\begin{poster}
    %%%%%%%%%%%%%%%%%%%%%%%%%%%%%%%%%%%%%%%%%
    %%             Begin poster            %%
    %%%%%%%%%%%%%%%%%%%%%%%%%%%%%%%%%%%%%%%%%
    
    \section{Introdu o}
            \lipsum[3]
        
    \section{Materiais e m todos}%
        
        \subsection{Coloque aqui seu modelo}
        
        \lipsum[11]
        
        \vspace{0.1cm}
        
        \begin{figure}
            \centering
            \captionsetup{type=figure}
            %mude o scale at  sua imagem ficar do tamanho certo o scale pode ter valores menores que 
            \includegraphics[scale=0.8]{images/exemplo.jpg}
            \caption{Trilha de desgaste topografia}
            \label{fig:lstm}
        \end{figure}
        
        \subsection{Coloque aqui seu \textit{dataset}}
        
            \lipsum[54]
            
            \vspace{0.5cm}
            %Aaqui come a a tabela
            \begin{table}
                \centering
                \captionsetup{type=table}
                \caption{\textit{Corpus} utilizados no estudo}
                \label{Corpus}
                \renewcommand{\arraystretch}{1.2}
                \resizebox{0.47\textwidth}{!}{%
                \begin{tabular}{lcccclcl}
                    \hline
                    &\textbf{Corpus}    &  \textbf{Caracteres  nicos} &    &  \textbf{Total de linhas}  &  & \textbf{Total de caracteres} &  \\ \hline
                    &SBSEThesis         & 88  &    &    2.311     &        & 771.179 &  \\ 
                    &Bible              & 63  &    &    32.359   &        & 3.924.374 &  \\
                    &JavaCode           & 69  &    &    436.565  &        & 12.053.424 &  \\ \hline
                \end{tabular}
            }
            \end{table}
            \vspace{0.5cm}
            
            \lipsum[57]
            
            \citep{defects2j}.
    
            \subsection{Subsection}
                \lipsum[8]
    
        \section{Resultados}%
                
            \lipsum[4]    
    
            \begin{figure}
                \centering
                \captionsetup{type=figure}
                \includegraphics[scale=1.3]{images/curva stribeck.jpg}
                \caption{Curva de stribeck \textit{}.}
                \label{fig:result}
            \end{figure}
    
            \lipsum[5]
            \cite{chollet2015keras}     
    
        \section{Considera es finais}
    
            \lipsum[15]
    
        \section{Agradecimento}
    
            \lipsum[57]
    
        
        
    
    %%%%%%%%%%%%%%%%%%%%%%%%%%%%%%%%%%%%%%%%%
    %%               End poster            %%
    %%%%%%%%%%%%%%%%%%%%%%%%%%%%%%%%%%%%%%%%%
    
    \end{poster}
\end{document}